\scanpage{21}
On peut aussi expliciter ainsi la construction de l'extension \(G\) \inserted{de \(H\) par \(N\)} associée à l'ext \(\deleted{\unsure{H}}G_0\) de \(\unsure{\mathfrak{S}|\mathcal{G}|\mathfrak{G}}\) par \(N\) et le \(\mathfrak{S}\)-hom \(\varphi:\deleted{\unsure{\mathfrak{S}}}\pi\to\mathfrak{z}(N)\) :

Soit \(G_{!}^{\deleted{\unsure{\xi|\varphi}}}\) l'ext. de \(H\) par \(N\) image inverse de l'ex \(G_0\) (de \(\mathfrak{S}\) par \(N\)) via \(H\to\mathfrak{S}\). Considérons l'extension \(\widetilde H\) de \(H\) par \(\pi\), \unsure{lé|le} \(\mathfrak{S}\)-hom \(\varphi:\pi\to\mathfrak{z}(N)\) nous donne une extn \unsure{corresp.|centrale} \(\widetilde H^{\varphi}_{\deleted{\unsure{\xi|\ast}}}\) de \(H\) par \(\mathfrak{z}(N)\), correspondant aux \(=\) opérations de \(H\) (et \(\mathfrak{S}\)) sur \(\mathfrak{z}(N)\) que l'ext. \(G_1\). On peut donc prendre le produit contracté des deux \inserted{extensions} \(G_1\) et \(\widetilde H^\varphi\), on trouve le \(G\) cherché.

D'ailleurs, \deleted{\unsure{soit}} comme \(G\) est un \unsure{revêtmt|revêtement} de \(H\), on aura
\[
 \pi_i(G)\xrightarrow{\sim}\pi_i(H)\quad\text{si }i\geq2
\]
\margin{\unsure{\(1\)}}
\diagram{\begin{tikzcd}[ampersand replacement=\&,column sep=.45em,row sep=.55em]
\pi_1(G)=\mathrm{Ker}\bigl( \& \pi_1(H)\arrow[r]\arrow[d,equal] \& \pi_0(N)\arrow[d,equal] \& \bigr)\\
\& \pi\arrow[r]\arrow[dr,"\varphi"'] \& N\arrow[d,phantom,"\rotatebox{90}{$\subset$}"] \&\\
\& \&\mathfrak{z}(N)\&
\end{tikzcd}}
\diagram{\begin{tikzcd}[ampersand replacement=\&,column sep=1em,row sep=.6em]
\pi_1(H)\arrow[r]\arrow[d,equal] \& N\arrow[r]\arrow[d,phantom,"\rotatebox{90}{$\subset$}"] \& \pi_0(G)\arrow[r]\&\pi_0(H)\arrow[r]\arrow[d,equal]\&1\\
\pi\arrow[r,"\varphi"']\&\mathfrak{z}(N)\&\&\mathfrak{S}\&
\end{tikzcd}}
\margin{i.e. \(\pi_0G\) est une \unsure{ext.|cat.} de \(\mathfrak{S}\) par \(\mathrm{Coker}(\varphi:\pi\to N)\)}

et on vérifie que l'on a une suite exacte
\[
 \deleted{\pi_0(\unsure{G})=}\qquad
 \pi\longrightarrow G_0\longrightarrow\pi_0(G)\longrightarrow1
\]
\[
 \deleted{\pi_0(G)\to G_0(\unsure{\mathfrak{S}})}
\]
(i.e. \(\pi_0(G)=\mathrm{Coker}(\pi\to G_0)\)) où

\(\pi\to G_0\) est le composé
\[
 \pi\longrightarrow\mathfrak{z}(N)\lhook\joinrel\longrightarrow N\lhook\joinrel\longrightarrow G_0\qquad \ldots\,..
\]

\textcircled{\(B'\)} Supposons seulement \(\pi,N,H\) discrets, \(\pi\subset\widetilde H_0\) (conditions \(\beta)\) et \(\gamma)\) mais pas néc. \(\alpha)\) (i.e. que \(\widetilde H_0\) soit connexe ou simplement connexe).

\(N\) étant discret, on a encore
\[
 \pi_i(G)\xrightarrow{\sim}\pi_i(H)\quad\text{si }i\geq2
\]
\hfill TSVP
% archive sheet number: 20

